\subsubsection{Experimentos Combinados sobre Fechas y Manejo de Cantidades}

En la \textit{notebook} 14\footnote{Disponible en: \url{https://github.com/aditzend/hyperpersonalization/blob/main/app/notebooks/14.ipynb}} se realiza un experimento que combina la extracción de información de conversaciones con el razonamiento cronológico y el manejo de datos personales utilizando GPT-3.5 y GPT-4. A continuación se detallan los pasos principales del experimento:

\paragraph{Carga de Librerías y Configuración}

Se utilizan FAISS para realizar búsquedas semánticas y OpenAI Embeddings para vectorizar los mensajes de las conversaciones. También se configura una función para interactuar con los modelos GPT-3.5 y GPT-4 a través de la API de OpenAI.

\begin{APACode}
import os
import requests
import json
from langchain.embeddings.openai import OpenAIEmbeddings
from langchain.vectorstores import FAISS

OPENAI_API_KEY = os.getenv("OPENAI_API_KEY")

index_name = "conversations"
index_path = "../indexes/conversations"
embeddings = OpenAIEmbeddings()

def gpt_system_user(
    system_message: str, user_message: str, model: str = "gpt-3.5-turbo"
):
    """Para usar en notebooks"""
    try:
        response = requests.post(
            url="https://api.openai.com/v1/chat/completions",
            params={"temperature": "0"},
            headers={
                "Content-Type": "application/json",
                "Authorization": f"Bearer {OPENAI_API_KEY}",
            },
            data=json.dumps({
                "model": model,
                "messages": [
                    {"content": system_message, "role": "system"},
                    {"content": user_message, "role": "user"},
                ]
            }),
            timeout=10,
        )
        choices = response.json()["choices"]
        answer = {
            "status_code": response.status_code,
            "text": choices[0]["message"],
        }
        return answer
    except requests.exceptions.RequestException:
        print("HTTP Request failed")
        return None
\end{APACode}

\paragraph{Simulación de Conversaciones}

Se crean contextos de conversación donde el usuario interactúa con el asistente virtual para realizar acciones como comprar pasajes para su familia de Buenos Aires a Córdoba o mejorar su categoría en un programa de viajero frecuente. Los datos proporcionados por el usuario, como nombres, edades y millas disponibles, se extraen y se almacenan en FAISS para su posterior consulta.

Ejemplos de mensajes:
\begin{itemize}
\item Usuario: ``Voy a sacar un pasaje para la familia de Buenos Aires a Córdoba''.
\item Asistente: ``¿Quiénes son los pasajeros?''.
\item Usuario: ``Ana es mi hija mayor, tiene 20 años. Juan es mi hijo menor, tiene 10 años. Mi esposa se llama María''.
\item Usuario: ``Quiero subir de categoría en el programa de viajero frecuente''.
\end{itemize}

\begin{APACode}
condensed_contexts = [
    {
        "messages": [
            {
                "role": "user",
                "content": "Voy a sacar un pasaje para la familia de buenos aires a cordoba"
            },
            {
                "role": "assistant",
                "content": "quienes son los pasajeros?"
            },
            {
                "role": "user",
                "content": "Ana es mi hija mayor, tiene 20 años. Juan es mi hijo menor, tiene 10 años. Mi esposa se llama María"
            },
            {
                "role": "assistant",
                "content": "en que fecha quiere viajar?"
            },
            {
                "role": "user",
                "content": "del 20 de agosto al 30 de agosto"
            }
        ],
        "metadata": {
            "session_id": "1",
            "person_id": "20",
            "context_id": "1",
            "created_at": "20200825T12:00:00.000Z",
        }
    }
]
\end{APACode}

\paragraph{Extracción de Datos}

El modelo GPT-3.5 se utiliza para extraer información relevante de las conversaciones, como los nombres, las edades, las millas acumuladas y las fechas de viaje. Esta información se almacena en FAISS junto con los metadatos correspondientes para facilitar su recuperación en futuras consultas.

Ejemplo de extracción: ``Pedro tiene una hija llamada Ana de 20 años y un hijo llamado Juan de 10 años. Quiere viajar con su familia de Buenos Aires a Córdoba del 20 al 30 de agosto''.

\begin{APACode}
db = FAISS.from_texts(texts=['notebook-14'], embedding=embeddings)
for context in condensed_contexts:
    messages = ''.join([str(message) for message in context['messages']])
    extraction = gpt_system_user(
        system_message="""
Extract every meaningful information about the user from the provided conversation.
Ask yourself, what happened in the conversation?
What data about the user did you learn?
At what date and time was this data gathered?

Create a summary of everything you learned about the user. Use this format:
On [date] at [time] the [user's / user's related person] [action] [data learned about the user].
        """,
        user_message=f"METADATA: ```{context['metadata']}``` CONVERSATION: ```{messages}```",
        model="gpt-3.5-turbo",
    )
    
    if extraction['status_code'] != 200:
        print(f"Error: OpenAI API returned status code {extraction['status_code']}")
        continue
    else:
        print(f"{extraction['text']['content']}")
        db.add_texts(
            texts=[extraction['text']['content']],
            metadatas=[context['metadata']],
        )
\end{APACode}

\paragraph{Consultas sobre Millas y Viajes}

Se plantea una pregunta como ``¿Cuántas millas tenía en el momento de hacer el viaje a Córdoba?''. El sistema intenta combinar la información obtenida de varios contextos para responder a la consulta. Sin embargo, se detecta un error en el razonamiento: el sistema responde que Pedro tenía 200 millas al momento del viaje, lo cual es incorrecto, ya que las millas se usaron para el pasaje.

\begin{APACode}
user_input = "cuantas millas tenia en el momento de hacer el viaje a cordoba?"
filter = {"person_id": "20"}
db_results = db.similarity_search_with_score(
    query=user_input, 
    embeddings=embeddings, 
    filter=filter, 
    k=2
)

context_k2 = [f'\n {db_results[i][0].page_content} {db_results[i][0].metadata}\n' 
              for i in range(len(db_results))]

system_prompt = f"""
You are the best personal assistant of SAIA Air. Your job is to help users in their travel needs.
Use the user's name in your answers.
You will find past conversations between you and the user between backticks.
Use them to answer the user's question.
To reason about dates, have in mind that today is September 13th, 2023.
Take a deep breath and work on this problem step-by-step.

```{context_k2}```
"""

answer = gpt_system_user(
    system_message=system_prompt, 
    user_message=user_input, 
    model="gpt-3.5-turbo"
)
\end{APACode}

\paragraph{Conclusiones de la catorceava iteración}

El sistema respondió: ``Según la información recopilada en las conversaciones anteriores, en el momento de hacer el viaje a Córdoba, Pedro tenía 200 millas disponibles''. 

Este resultado presenta un error de razonamiento, ya que no considera correctamente el orden cronológico de los eventos ni el uso de las millas para el pasaje.

\paragraph{Posibles Mejoras}

Se identifican algunos problemas de razonamiento, particularmente al intentar unir información de diferentes contextos. Se proponen futuros pasos para mejorar el modelo:

\begin{itemize}
\item Navegar el grafo de relaciones del usuario y cambiar el foco de la conversación (``este turno es para mi madre'')
\item Usar el mismo método para guardar información como ``La madre del usuario tiene un turno para el dentista el 1 de enero de 2021''
\item Agregar más datos personales para evitar confusiones en el futuro
\item Simular llamados a APIs desde diferentes relaciones del usuario y evaluar posibles confusiones
\end{itemize}

La \textit{notebook} 14 muestra un experimento que combina la extracción de datos y el razonamiento sobre múltiples contextos. Si bien el sistema es capaz de recuperar y manejar información relevante, presenta errores al unir datos de diferentes contextos, especialmente cuando se trata de razonamiento cronológico y manejo de cantidades, como las millas disponibles. 